                                                                                              n
Problem 6. For a permutation σ = (i1 , i2 , ..., in ) of (1, 2, ..., n) define D(σ) =
                                                                                              P
                                                                                                    |ik − k|. Let Q(n, d) be                             IMC2008, Blagoevgrad, Bulgaria
                                                                                              k=1
the number of permutations σ of (1, 2, ..., n) with d = D(σ). Prove that Q(n, d) is even for d ≥ 2n.                                                                   Day 1, July 27, 2008
Solution. Consider the n × n determinant

                                                       1           x    . . . xn−1
                                                       x           1    . . . xn−2                                             Problem 1. Find all continuous functions f : R → R such that f (x) − f (y) is rational for all reals x and
                                         ∆(x) =        ..          ..   ..      ..
                                                        .           .       .    .                                             y such that x − y is rational.
                                                     xn−1 xn−2          ...    1                                               Solution. We prove that f (x) = ax + b where a ∈ Q and b ∈ R. These functions obviously satify the
                                                                                                                               conditions.
where the ij-th entry is x|i−j| . From the definition of the determinant we get                                                    Suppose that a function f (x) fulfills the required properties. For an arbitrary rational q, consider the
                                              X                                                                                function gq (x) = f (x+ q) −f (x). This is a continuous function which attains only rational values, therefore
                                   ∆(x) =            (−1)inv(i1 ,...,in) xD(i1 ,...,in)                                        gq is constant.
                                             (i1 ,...,in)∈Sn
                                                                                                                                   Set a = f (1) − f (0) and b = f (0). Let n be an arbitrary positive integer and let r = f (1/n) − f (0).
                                                                                                                               Since f (x + 1/n) − f (x) = f (1/n) − f (0) = r for all x, we have
where Sn is the set of all permutations of (1, 2, ..., n) and inv(i1 , ..., in ) denotes the number of inversions in
the sequence (i1 , ..., in ). So Q(n, d) has the same parity as the coefficient of xd in ∆(x).                                        f (k/n) − f (0) = (f (1/n) − f (0)) + (f (2/n) − f (1/n)) + . . . + (f (k/n) − f ((k − 1)/n) = kr
   It remains to evaluate ∆(x). In order to eliminate the entries below the diagonal, subtract the (n−1)-th
row, multiplied by x, from the n-th row. Then subtract the (n − 2)-th row, multiplied by x, from the                           and
(n − 1)-th and so on. Finally, subtract the first row, multiplied by x, from the second row.
                                                                                                                               f (−k/n) − f (0) = −(f (0) − f (−1/n)) − (f (−1/n) − f (−2/n)) − . . . − (f (−(k − 1)/n) − f (−k/n) = −kr
             1      x     . . . xn−2 xn−1           1   x                ...      xn−2       xn−1
             x      1     . . . xn−3 xn−2           0 1 − x2             . . . xn−3 − xn−1 xn−2 − xn                           for k ≥ 1. In the case k = n we get a = f (1) − f (0) = nr, so r = a/n. Hence, f (k/n) − f (0) = kr = ak/n
 ∆(x) =      ..     ..    ..      ..   .. = . . . = ..   ..              ..         ..         ..    = (1 − x2 )n−1 .          and then f (k/n) = a · k/n + b for all integers k and n > 0.
              .      .        .    .    .            .    .                  .       .          .
           xn−2 xn−3      ...    1    x             0   0                ...         1 − x2    x − x3                              So, we have f (x) = ax + b for all rational x. Since the function f is continous and the rational numbers
           xn−1 xn−2      ...    x    1             0   0                ...           0       1 − x2                          form a dense subset of R, the same holds for all real x.

   For d ≥ 2n, the coefficient of xd is 0 so Q(n, d) is even.                                                                  Problem 2. Denote by V the real vector space of all real polynomials in one variable, and let P : V → R
                                                                                                                               be a linear map. Suppose that for all f, g ∈ V with P (f g) = 0 we have P (f ) = 0 or P (g) = 0. Prove that
                                                                                                                               there exist real numbers x0 , c such that P (f ) = c f (x0 ) for all f ∈ V .
                                                                                                                               Solution. We can assume that P 6= 0.
                                                                                                                                   Let f ∈ V be such that P (f ) 6= 0. Then P (f 2 ) 6= 0, and therefore P (f 2) = aP (f ) for some non-zero
                                                                                                                               real a. Then 0 = P (f 2 − af ) = P (f (f − a)) implies P (f − a) = 0, so we get P (a) 6= 0. By rescaling, we
                                                                                                                               can assume that P (1) = 1. Now P (X + b) = 0 for b = −P (X). Replacing P by P̂ given as

                                                                                                                                                                         P̂ (f (X)) = P (f (X + b))

                                                                                                                               we can assume that P (X) = 0.
                                                                                                                                  Now we are going to prove that P (X k ) = 0 for all k ≥ 1. Suppose this is true for all k < n. We know
                                                                                                                               that P (X n + e) = 0 for e = −P (X n ). From the induction hypothesis we get

                                                                                                                                                               P (X + e)(X + 1)n−1 = P (X n + e) = 0,
                                                                                                                                                                                     


                                                                                                                               and therefore P (X + e) = 0 (since P (X + 1) = 1 6= 0). Hence e = 0 and P (X n ) = 0, which completes the
                                                                                                                               inductive step. From P (1) = 1 and P (X k ) = 0 for k ≥ 1 we immediately get P (f ) = f (0) for all f ∈ V .




                                                               4                                                                                                                     1
Problem 3. Let p be a polynomial with integer coefficients and let a1 < a2 < . . . < ak be integers.                     is a set of three special triples also (we may suppose that a + b + c < 1, because otherwise all three triples
                                                                                                                         are equal and our statement is trivial).
  a) Prove that there exists a ∈ Z such that p(ai ) divides p(a) for all i = 1, 2, . . . , k.
                                                                                                                             If there is a special triple (x, y, z) which is not worse than any triple from S1 , then the triple
  b) Does there exist an a ∈ Z such that the product p(a1 ) · p(a2 ) · . . . · p(ak ) divides p(a)?
                                                                                                                                               ((1 − a − b − c)x + a, (1 − a − b − c)y + b, (1 − a − b − c)z + c)
Solution. The theorem is obvious if p(ai ) = 0 for some i, so assume that all p(ai ) are nonzero and pairwise
different.                                                                                                               is special and not worse than any triple from S. We also have a(S1 ) = b(S1 ) = c(S1 ) = 0, so we may
    There exist numbers s, t such that s|p(a1 ), t|p(a2 ), st = lcm(p(a1 ), p(a2 )) and gcd(s, t) = 1.                   suppose that the same holds for our starting set S.
    As s, t are relatively prime numbers, there exist m, n ∈ Z such that a1 + sn = a2 + tm =: b2 . Obviously                  Suppose that one element of S has two entries equal to 0.
s|p(a1 + sn) − p(a1 ) and t|p(a2 + tm) − p(a2 ), so st|p(b2 ).                                                                Note that one of the two remaining triples from S is not worse than the other. This triple is also not
    Similarly one obtains b3 such that p(a3 )|p(b3 ) and p(b2 )|p(b3 ) thus also p(a1 )|p(b3 ) and p(a2 )|p(b3 ).        worse than all triples from S because any special triple is not worse than itself and the triple with two
    Reasoning inductively we obtain the existence of a = bk as required.                                                 zeroes.
    The polynomial p(x) = 2x2 + 2 shows that the second part of the problem is not true, as p(0) = 2,                         So we have a = b = c = 0 but we may suppose that all triples from S contain at most one zero. By
p(1) = 4 but no value of p(a) is divisible by 8 for integer a.                                                           transposing triples and elements in triples (elements in all triples must be transposed simultaneously) we
Remark. One can assume that the p(ai ) are nonzero and ask for a such that p(a) is a nonzero mul-                        may achieve the following situation x1 = y2 = z3 = 0 and x2 ⩾ x3 . If z2 ⩾ z1 , then the second triple
tiple of all p(ai ). In the solution above, it can happen that p(a) = 0. But every number p(a +                          (x2 , 0, z2 ) is not worse than the other two triples from S. So we may assume that z1 ⩾ z2 . If y1 ⩾ y3 ,
np(a1 )p(a2 ) . . . p(ak )) is also divisible by every p(ai ), since the polynomial is nonzero, there exists n such      then the first triple is not worse than the second and the third and we assume y3 ⩾ y1 . Consider the
that p(a + np(a1 )p(a2 ) . . . p(ak )) satisfies the modified thesis.                                                    three pairs of numbers x2 , y1; z1 , x3 ; y3 , z2 . The sum of all these numbers is three and consequently the
                                                                                                                         sum of the numbers in one of the pairs is less than or equal to one. If it is the first pair then the triple
Problem 4. We say a triple (a1 , a2 , a3 ) of nonnegative reals is better than another triple (b1 , b2 , b3 ) if two     (x2 , 1 − x2 , 0) is not worse than all triples from S, for the second we may take (1 − z1 , 0, z1 ) and for the
out of the three following inequalities a1 > b1 , a2 > b2 , a3 > b3 are satisfied. We call a triple (x, y, z)            third — (0, y3, 1 − y3 ). So we found a desirable special triple for any given S.
special if x, y, z are nonnegative and x + y + z = 1. Find all natural numbers n for which there is a set S
of n special triples such that for any given special triple we can find at least one better triple in S.                 Problem 5. Does there exist a finite group G with a normal subgroup H such that |Aut H| > |Aut G|?
Solution. The answer is n ⩾ 4.                                                                                           Solution. Yes. Let H be the commutative group H = F32 , where F2 ∼
                                                                                                                                                                                          = Z/2Z is the field with two elements.
    Consider the following set of special triples:                                                                       The group of automorphisms of H is the general linear group GL3 F2 ; it has
                                                                                
                            8 7            2      3         3 2               2 11 2
                          0, ,      ,        , 0,     ,      , ,0 ,            , ,       .                                                                 (8 − 1) · (8 − 2) · (8 − 4) = 7 · 6 · 4 = 168
                            15 15          5      5         5 5              15 15 15
We will prove that any special triple (x, y, z) is worse than one of these (triple a is worse than triple b if           elements. One of them is the shift operator φ : (x1 , x2 , x3 ) 7→ (x2 , x3 , x1 ).
triple b is better than triple a). We suppose that some special triple (x, y, z) is actually not worse than the             Now let T = {a0 , a1 , a2 } be a group of order 3 (written multiplicatively); it acts on H by τ (a) = φ. Let
first three of the triples from the given set, derive some conditions on x, y, z and prove that, under these             G be the semidirect product G = H ⋊τ T . In other words, G is the group of 24 elements
conditions, (x, y, z) is worse than the fourth triple    from the set.
                                           8 7                         8          7                                                                     G = {bai : b ∈ H, i ∈ (Z/3Z)},        ab = φ(b)a.
                                                  
    Triple (x,y, z) is not worse than 0, 15 , 15 means that y ⩾ 15      or z ⩾ 15  . Triple (x, y, z) is not worse
than 52 , 0, 35 — x ⩾ 25 or z ⩾ 53 . Triple (x, y, z) is not worse than 35 , 52 , 0 — x ⩾ 35 or y ⩾ 52 . Since
                                                                                                                         G has one element e of order 1 and seven elements b, b ∈ H, b 6= e of order 2.
x + y + z = 1, then it is impossible that all inequalities x ⩾ 25 , y ⩾ 25 and z ⩾ 15    7
                                                                                           are true. Suppose that
                                                                                                                            If g = ba, we find that g 2 = baba = bφ(b)a2 6= e, and that
      2
x < 5 , then y ⩾ 5 and z ⩾ 5 . Using x + y + z = 1 and x ⩾ 0 we get x = 0, y = 52 , z = 35 . We obtain
                      2         3

the triple 0, 52 , 53 which is worse than 152 11 2
                                              , 15 , 15 . Suppose that y < 25 , then x ⩾ 53 and z ⩾ 15     7
                                                                                                              and this                              g 3 = bφ(b)a2 ba = bφ(b)aφ(b)a2 = bφ(b)φ2 (b)a3 = ψ(b),
                                                                            7               2              8
is a contradiction to the admissibility of (x, y, z). Suppose that z < 15 , then x ⩾ 5 and y ⩾ 15           . We get
                                     1
(by admissibility, again) that z ⩽ 15  and y ⩽ 35 . The last inequalities imply that 15    2 11 2
                                                                                             , 15 , 15 is better than    where the homomorphism ψ : H → H is defined as ψ : (x1 , x2 , x3 ) 7→ (x1 + x2 + x3 )(1, 1, 1). It is clear that
(x, y, z).                                                                                                               g 3 = ψ(b) = e for 4 elements b ∈ H, while g 6 = ψ 2 (b) = e for all b ∈ H.
    We will prove that for any given set of three special triples one can find a special triple which is not                 We see that G has 8 elements of order 3, namely ba and ba2 with b ∈ Ker ψ, and 8 elements of order 6,
worse than any triple from the set. Suppose we have a set S of three special triples                                     namely ba and ba2 with b 6∈ Ker ψ. That accounts for orders of all elements of G.
                                   (x1 , y1 , z1 ),   (x2 , y2 , z2 ),   (x3 , y3, z3 ).                                     Let b0 ∈ H \Kerψ be arbitrary; it is easy to see that G is generated by b0 and a. As every automorphism
                                                                                                                         of G is fully determined by its action on b0 and a, it follows that G has no more than
Denote a(S) = min(x1 , x2 , x3 ), b(S) = min(y1 , y2 , y3), c(S) = min(z1 , z2 , z3 ). It is easy to check that S1 :
                                                                                                                                                                            7 · 8 = 56
                                                                                  
                                      x1 − a           y1 − b        z1 − c
                                               ,                  ,
                                   1−a−b−c 1−a−b−c 1−a−b−c
                                                                                                                       automorphisms.
                                      x2 − a           y2 − b        z2 − c
                                               ,                  ,                                                      Remark. G and H can be equivalently presented as subgroups of S6 , namely as H = h(12), (34), (56)i and
                                   1−a−b−c 1−a−b−c 1−a−b−c
                                                                                                                       G = h(135)(246), (12)i.
                                      x3 − a           y3 − b        z3 − c
                                               ,                  ,
                                   1−a−b−c 1−a−b−c 1−a−b−c

                                                             2                                                                                                                  3
